\documentclass{article}
\usepackage{amsmath,amssymb,amsthm}
\newtheorem{theorem}{Theorem}
\newtheorem{corollary}{Corollary}
\newcommand{\TT}{\mathrm{TT}}
\newcommand{\rank}{\operatorname{rank}}

\begin{document}
\section{Particle-TT obstruction for alternating tensors}

Let $V$ be a finite-dimensional real or complex Hilbert space and identify
$\Lambda^N V$ with the alternating subspace of $V^{\otimes N}$. For
$C\in V^{\otimes N}$, write $C_{(k)}$ for its $k\mid N-k$ particle
matricization.

\begin{theorem}[Exact particle-TT ranks]
For $1\le k<N$, the minimal exact ordinary-TT bond dimension satisfies
\[
 r_k^{\TT}(C)=\rank C_{(k)}.
\]
\end{theorem}
\begin{proof}
Any TT cut factors $C_{(k)}$ through its bond space, which gives the lower
bound. Successive exact singular-value decompositions attain all unfolding
ranks simultaneously.
\end{proof}

\begin{theorem}[Exchange rank floor]
If $0\ne C\in\Lambda^N V$, then
\[
 \rank C_{(k)}\ge \binom{N}{k},\qquad 1\le k<N.
\]
The bound is attained by every nonzero decomposable $N$-form.
\end{theorem}
\begin{proof}
Choose a nonzero alternating coefficient $c_I$ supported on an $N$-element
index set $I$. Restrict rows to $k$-subsets $A\subset I$ and match them to
columns $I\setminus A$. The resulting $\binom{N}{k}$ square submatrix is
diagonal up to signs, with diagonal entries $\pm c_I$: an off-diagonal entry
repeats an index and vanishes. A decomposable form supported on an
$N$-dimensional space has no larger contraction image.
\end{proof}

\begin{corollary}[No strict-antisymmetry truncation below the floor]
If $\widetilde C\ne0$ and
$r_k^{\TT}(\widetilde C)<\binom{N}{k}$ for some $k$, then
$\widetilde C\notin\Lambda^N V$.
\end{corollary}

Use the normalized Slater convention
\[
 \Phi=\frac1{\sqrt{N!}}\sum_{\pi\in S_N}\mathrm{sgn}(\pi)
 u_{\pi(1)}\otimes\cdots\otimes u_{\pi(N)},
\]
where the $u_i$ are orthonormal.

\begin{theorem}[Flat particle-Schmidt spectrum]
At the $k\mid N-k$ particle cut,
\[
 \Phi=\frac1{\sqrt{\binom{N}{k}}}\sum_{|I|=k}
 (-1)^{\epsilon(I)}\Phi_I\otimes\Phi_{I^c}.
\]
This is a Schmidt decomposition. Its $\binom{N}{k}$ nonzero singular values
all equal $\binom{N}{k}^{-1/2}$.
\end{theorem}
\begin{proof}
Partition the signed permutation sum by the orbital subset occupying the first
$k$ particle slots. Each class factors into normalized internal Slater sums.
Its coefficient is
$\sqrt{k!(N-k)!/N!}=\binom{N}{k}^{-1/2}$. Distinct orbital subsets give
orthonormal Slater states on each side.
\end{proof}

\begin{corollary}[Approximate bond lower bound]
For $0\le r\le \binom{N}{k}$, the best cut-rank-$r$ approximation has relative
squared error
\[
 \epsilon_{k,r}^2=1-\frac{r}{\binom{N}{k}}.
\]
Thus, for $0\le\epsilon\le1$, every ordinary TT with relative error at most
$\epsilon$ must obey
\[
 r_k^{\TT}\ge\left\lceil(1-\epsilon^2)\binom{N}{k}\right\rceil.
\]
\end{corollary}
\begin{proof}
Apply the Eckart--Young theorem to the flat singular-value list.
\end{proof}

These statements concern ordinary particle-site TT coordinates. They do not
exclude ordered-sector, exterior-carrier, determinant-sum, or occupation-space
representations. Nor is spectral flatness asserted for a general element of
$\Lambda^N V$.
\end{document}
